% GENERATED FILE. Edit canonical lessons, then run npm run generate:books.
% canonical-source-hash: fnv1a64:d35024b67a95184d
% canonical-lessons: MR-C278-vacat-hoto, MR-C278-astana, MR-C278-jayco, MR-R278-the-past-three-ways

\chapter{Was Doing, While, Used To}
\label{ch:mr-was-doing-while-used-to}
\hlchaptermodality{278}

\begin{chapteropening}
\textbf{By the end of this chapter:} \emph{I can set the scene with vācat hoto and astānā, and say what used to happen with jāyco.}

Complete MR-R278-the-past-three-ways: tell what happened, what was going on and what used to happen.
\end{chapteropening}

\section[\texorpdfstring{\mr{मी} \mr{वाचत} \mr{होतो} (mī vācat hoto)}{मी वाचत होतो (mī vācat hoto)}]{\mr{मी} \mr{वाचत} \mr{होतो} (mī vācat hoto) — what was going on}

\label{lesson:MR-C278-vacat-hoto}



\begin{quote}
Say \textquotedblleft{}I won't come tomorrow\textquotedblright{}. (\emph{mī udyā yeṇār nāhī}.) Then \textquotedblleft{}I was at home\textquotedblright{} and \textquotedblleft{}I read a book\textquotedblright{}. (\emph{mī gharī hoto}, \emph{mī pustak vāclaṁ}.)

The last is an event: it happened and it was over. A story also needs what was going on at the time.
\end{quote}

\subsection*{You'll want to know: -\mr{त} \mr{होतो}, 'was doing'}
\begin{quote}
\textbf{\mr{काल} \mr{रात्री} \mr{नऊ} \mr{वाजता} \mr{मी} \mr{पुस्तक} \mr{वाचत} \mr{होतो}.}
\end{quote}

\begin{quote}
\emph{kāl rātrī naū vājtā mī pustak vācat hoto.} — \textquotedblleft{}At nine last night I was reading a book.\textquotedblright{}
\end{quote}

Two pieces. \textbf{\mr{वाचत}} is the bare \textbf{-\mr{त}} shape you know from \textbf{\mr{आवडत} \mr{नाही}}; \textbf{\mr{होतो}}, \textquotedblleft{}I was\textquotedblright{}, carries the person and the gender:

\begin{itemize}
\raggedright
  \item \textbf{\mr{ती}} \textbf{\mr{वाचत} \mr{होती}} (\emph{vācat hotī}) — she was reading
  \item \textbf{\mr{तो}} \textbf{\mr{काम} \mr{करत} \mr{होता}} (\emph{kām karat hotā}) — he was working
  \item \textbf{\mr{आम्ही}} \textbf{\mr{बोलत} \mr{होतो}} (\emph{bolat hoto}) — we were talking
  \item \textbf{\mr{पाऊस} \mr{पडत} \mr{होता}} (\emph{pāūs paḍat hotā}) — it was raining
\end{itemize}

This past follows the doer even with an object: \textbf{\mr{ती} \mr{पुस्तक} \mr{वाचत} \mr{होती}}. No \textbf{-\mr{ने}} here.

\begin{grammarlens}[title={Grammar Lens: the scene and the event}]
\begin{quote}
\textbf{\mr{मी} \mr{वाचत} \mr{होतो}. \mr{तेव्हा} \mr{मीरा} \mr{आली}.}
\end{quote}

\begin{quote}
\emph{mī vācat hoto. tevhā mīrā ālī.} — \textquotedblleft{}I was reading. At that moment, Meera came.\textquotedblright{}
\end{quote}

The first sentence is the \textbf{scene}: it was going on. \textbf{\mr{आली}} is the \textbf{event}: it happened once and cut into the scene. \textbf{\mr{तेव्हा}}, the \textquotedblleft{}then\textquotedblright{} of \textbf{\mr{जेव्हा} … \mr{तेव्हा}}, hooks one onto the other. A told story moves on this rhythm.
\end{grammarlens}

\subsection*{Your turn}
Say these aloud:

\begin{itemize}
\raggedright
  \item \textquotedblleft{}She was reading\textquotedblright{} — \emph{tī vācat hotī}
  \item \textquotedblleft{}It was raining\textquotedblright{} — \emph{pāūs paḍat hotā}
  \item \textquotedblleft{}I was reading. At that moment, Meera came.\textquotedblright{} — \emph{mī vācat hoto. tevhā mīrā ālī.}
\end{itemize}

\begin{tcolorbox}[breakable,colback=teal!4,colframe=teal!35!black,title={Before you move on}]
\textquotedblleft{}I was reading\textquotedblright{}? (\textbf{\mr{मी} \mr{वाचत} \mr{होतो}}.) \textquotedblleft{}It was raining\textquotedblright{}? (\textbf{\mr{पाऊस} \mr{पडत} \mr{होता}}.) Which past is the scene? (\textbf{-\mr{त} \mr{होतो}}.) Which is the event? (\textbf{The plain past}: \textbf{\mr{आली}}.)
\end{tcolorbox}
\section[\texorpdfstring{\mr{असताना} (astānā)}{असताना (astānā)}]{\mr{असताना} (astānā) — while it was going on}

\label{lesson:MR-C278-astana}



\begin{quote}
Say \textquotedblleft{}I was reading\textquotedblright{}. (\emph{mī vācat hoto}.) Then the two-part \textquotedblleft{}when … then\textquotedblright{} frame. (\emph{jevhā … tevhā}.)
\end{quote}

\subsection*{You'll want to know: \mr{असताना}, 'while'}
\begin{quote}
\textbf{\mr{मी} \mr{वाचत} \mr{असताना} \mr{मीरा} \mr{आली}.}
\end{quote}

\begin{quote}
\emph{mī vācat astānā mīrā ālī.} — \textquotedblleft{}While I was reading, Meera came.\textquotedblright{}
\end{quote}

\textbf{\mr{असताना}} is \textbf{\mr{असणे}}, \textquotedblleft{}to be\textquotedblright{}, with \textbf{-\mr{ताना}}, \textquotedblleft{}while being\textquotedblright{}. Put it after the \textbf{-\mr{त}} shape, and the scene needs no second clause: \textbf{\mr{मी} \mr{वाचत} \mr{असताना}}. It follows a place or a state too: \textbf{\mr{मी} \mr{घरी} \mr{असताना}} (\emph{mī gharī astānā}), \textquotedblleft{}while I was at home\textquotedblright{}; \textbf{\mr{लहान} \mr{असताना}} (\emph{lahān astānā}), \textquotedblleft{}when (I was) small\textquotedblright{}.

It does not change for person or gender. The event after it does: \textbf{\mr{मीरा} \mr{आली}}.

\subsection*{You'll want to know: \mr{जेव्हा} … \mr{तेव्हा} with the past}
\begin{quote}
\textbf{\mr{मी} \mr{घरी} \mr{आलो}, \mr{तेव्हा} \mr{आई} \mr{चहा} \mr{करत} \mr{होती}.}
\end{quote}

\begin{quote}
\emph{mī gharī ālo, tevhā āī cahā karat hotī.} — \textquotedblleft{}When I came home, Mother was making tea.\textquotedblright{}
\end{quote}

\textbf{\mr{जेव्हा}} can stay or go: with the past, the first clause alone sets the time, and \textbf{\mr{तेव्हा}} picks it up.

\subsection*{Your turn}
Say these aloud:

\begin{itemize}
\raggedright
  \item \textquotedblleft{}While I was reading, Meera came\textquotedblright{} — \emph{mī vācat astānā mīrā ālī}
  \item \textquotedblleft{}While I was at home, it was raining\textquotedblright{} — \emph{mī gharī astānā pāūs paḍat hotā}
  \item \textquotedblleft{}When I came home, Mother was making tea\textquotedblright{} — \emph{mī gharī ālo, tevhā āī cahā karat hotī}
\end{itemize}

\begin{tcolorbox}[breakable,colback=teal!4,colframe=teal!35!black,title={Before you move on}]
\textquotedblleft{}While I was reading\textquotedblright{}? (\textbf{\mr{मी} \mr{वाचत} \mr{असताना}}.) \textquotedblleft{}When (I was) small\textquotedblright{}? (\textbf{\mr{लहान} \mr{असताना}}.) Does \textbf{\mr{असताना}} change for \textquotedblleft{}she\textquotedblright{}? (No.)
\end{tcolorbox}
\section[\texorpdfstring{\mr{मी} \mr{जायचो} (mī jāyco)}{मी जायचो (mī jāyco)}]{\mr{मी} \mr{जायचो} (mī jāyco) — what used to happen}

\label{lesson:MR-C278-jayco}



\begin{quote}
Say \textquotedblleft{}when (I was) small\textquotedblright{}. (\emph{lahān astānā}.) Then \textquotedblleft{}formerly\textquotedblright{}. (\emph{pūrvī}.)
\end{quote}

\subsection*{You'll want to know: -\mr{ायचो}, 'used to'}
\begin{quote}
\textbf{\mr{लहान} \mr{असताना} \mr{मी} \mr{रोज} \mr{नदीवर} \mr{जायचो}.}
\end{quote}

\begin{quote}
\emph{lahān astānā mī roj nadīvar jāyco.} — \textquotedblleft{}When I was small, I used to go to the river every day.\textquotedblright{}
\end{quote}

Take \textbf{-\mr{णे}} off and put \textbf{-\mr{ायच}-} on, then the endings of \textbf{\mr{गेलो}}:

\begin{itemize}
\raggedright
  \item \textbf{\mr{मी}} \textbf{\mr{जायचो}} (\emph{jāyco}) — I used to go; a woman says \textbf{\mr{जायचे}} (\emph{jāyce})
  \item \textbf{\mr{तो}} \textbf{\mr{जायचा}} (\emph{jāycā}), \textbf{\mr{ती}} \textbf{\mr{जायची}} (\emph{jāycī}) — he, she used to go
  \item \textbf{\mr{आम्ही}} \textbf{\mr{जायचो}}, \textbf{\mr{ते}} \textbf{\mr{जायचे}} — we, they used to go
\end{itemize}

\textbf{\mr{करणे}} $\to$ \textbf{\mr{करायचो}} (\emph{karāyco}); \textbf{\mr{राहणे}} $\to$ \textbf{\mr{राहायचो}} (\emph{rāhāyco}): \textbf{\mr{पूर्वी} \mr{मी} \mr{पुण्यात} \mr{राहायचो}} (\emph{pūrvī mī puṇyāt rāhāyco}), \textquotedblleft{}I used to live in Pune\textquotedblright{}. With an object it follows the doer: \textbf{\mr{मी} \mr{पुस्तक} \mr{वाचायचो}} (\emph{mī pustak vācāyco}).

\subsection*{Your turn}
Say these aloud:

\begin{itemize}
\raggedright
  \item \textquotedblleft{}When I was small, I used to go to the river every day\textquotedblright{} — \emph{lahān astānā mī roj nadīvar jāyco}
  \item \textquotedblleft{}Formerly she used to live in Pune\textquotedblright{} — \emph{pūrvī tī puṇyāt rāhāycī}
  \item \textquotedblleft{}We used to read books\textquotedblright{} — \emph{āmhī pustakaṁ vācāyco}
\end{itemize}

\textbf{\mr{पुस्तकं}} (\emph{pustakaṁ}, \textquotedblleft{}books\textquotedblright{}).

\begin{tcolorbox}[breakable,colback=teal!4,colframe=teal!35!black,title={Before you move on}]
\textquotedblleft{}I used to go\textquotedblright{}? (\textbf{\mr{जायचो}}, or \textbf{\mr{जायचे}}.) \textquotedblleft{}She used to live\textquotedblright{}? (\textbf{\mr{राहायची}}.) Which past is a habit? (\textbf{-\mr{ायचो}}.)
\end{tcolorbox}
\section[(review)]{Review — the past three ways}

\label{lesson:MR-R278-the-past-three-ways}



\begin{quote}
Say \textquotedblleft{}I am going to go\textquotedblright{} and \textquotedblleft{}I won't go\textquotedblright{}. (\emph{mī jāṇār āhe}, \emph{mī jāṇār nāhī}.) Then \textquotedblleft{}he was not at home\textquotedblright{} and \textquotedblleft{}she was reading\textquotedblright{}. (\emph{to gharī navhtā}, \emph{tī vācat hotī}.)
\end{quote}

\subsection*{Your turn}
Say these aloud:

\begin{itemize}
\raggedright
  \item At nine last night I was reading a book.
  \item While I was reading, Meera came.
  \item When I came home, Mother was making tea.
  \item When I was small, I used to go to the river every day.
  \item one thing you were doing yesterday evening, and one thing that cut into it
\end{itemize}

\begin{tcolorbox}[breakable,colback=teal!4,colframe=teal!35!black,title={Before you move on}]
Which past is the event? (\textbf{The plain past}: \textbf{\mr{आली}}.) Which is the scene? (\textbf{-\mr{त} \mr{होतो}}, or \textbf{\mr{असताना}}.) Which is the habit? (\textbf{-\mr{ायचो}}.) \textquotedblleft{}When I came home\textquotedblright{}? (\textbf{\mr{मी} \mr{घरी} \mr{आलो}, \mr{तेव्हा} …})
\end{tcolorbox}
